\section{Safety, Privacy, and Governance}
\label{sec:safety_privacy_governance}

Speech Large Language Models (Speech LLMs) extend language-model interaction from text to spoken communication, but this shift also changes the safety and privacy problem. Speech is not simply text represented as audio. It can serve as a command channel, an acoustic signal, a marker of speaker identity, and a carrier of synthetic media. A spoken utterance may encode words, speaker characteristics, prosody, emotion, accent, background sounds, and signal-level artifacts. These properties make speech useful for natural interaction, but they also create attack and misuse pathways that are not captured by text-only safety evaluation.

This section reviews six groups of risks and mitigations: input-side audio jailbreaks and adversarial speech, audio safety benchmarks, guardrails and anti-cloning defenses, synthetic voice misuse, speaker privacy protection, and documentation for governance. Because many of the cited studies are recent preprints, their findings should be interpreted cautiously. Still, the literature points to a consistent pattern: safeguards developed for text-only LLMs do not necessarily transfer to speech-capable systems. Speech LLMs therefore require safety analysis that accounts for audio representations, speech generation, speaker identity, and deployment context.

\subsection{Speech-Specific Risk Surface}

Speech LLMs expose semantic, acoustic, biometric, and synthetic-media risks at the same time. In text-only LLMs, harmful prompts are usually represented as discrete text tokens. In Speech LLMs, the input may pass through raw waveform encoders, speech tokenizers, ASR modules, acoustic projectors, or multimodal alignment layers before reaching the language backbone. These processing stages can change model behavior and create attack opportunities that are difficult to detect from transcripts alone. Recent AI-ethics work argues that audio language model deployments should be scoped by a principle of least privilege---granting a system access only to the acoustic information a task actually requires---precisely because raw audio carries far more than the intended semantic payload \citep{he2025modelhears}.

\subsubsection{Input-Side Attacks: Harmful Spoken Prompts and Modality Transfer Failure}

Red-teaming studies show that safety behavior can change when harmful requests are presented as speech rather than text. \textbf{Audio Is the Achilles' Heel} evaluates Qwen-Audio, Qwen2-Audio, SALMONN-7B, SALMONN-13B, and Gemini-1.5-Pro under three settings: harmful questions in text and audio formats, harmful text paired with distracting non-speech audio, and speech-specific jailbreaks. Beyond aggregate attack success, the paper analyzes representation-space behavior using t-SNE visualizations and response-opening patterns. The authors find that harmful questions can become more effective when moved into the audio modality for several open-source audio LMMs, and that non-speech audio can reshape representation spaces in ways that affect safety behavior \citep{yang2024audioachillesheelred}. This result suggests that safety alignment learned through text may weaken after multimodal adaptation.

Adversarial audio attacks further show that the waveform itself can be optimized to bypass safety constraints. \textbf{AdvWave} treats jailbreak attacks against large audio-language models as a waveform-level optimization problem. The method includes dual-phase optimization to address gradient shattering from audio encoder discretization, adaptive target search to handle response variability, and classifier-guided natural-sound constraints to keep perturbations perceptually plausible \citep{kang2024advwavestealthyadversarialjailbreak}. This matters because the attack can operate below the transcript level: the spoken request may remain understandable to humans, while the acoustic perturbation changes how the model processes or answers it.

Multimodal attacks add another path. \textbf{JAMA} proposes a joint audio-text attack against spoken language models by combining gradient-based text suffix optimization with audio perturbation. The method jointly optimizes across modalities and evaluates both fixed-prompt and prompt-adjusted settings. Experiments show that combined text-audio attacks can outperform unimodal attacks across several spoken language models \citep{krishnan2026optimizingmultimodaljailbreaksspoken}.

Sparse token attacks show that only parts of the audio sequence may be sufficient for a successful jailbreak. \textbf{TAGO} analyzes token-aligned gradients and finds that gradient energy is not evenly distributed across audio tokens. Its method identifies high-contribution token regions and optimizes only the corresponding waveform segments. Comparisons with full-waveform attacks, token-selection strategies, and efficiency analyses show that sparse token optimization can reduce attack cost while maintaining strong jailbreak performance \citep{fang2026sparsetokenssufficejailbreaking}. This suggests that small local changes in audio should not be treated as automatically harmless.

Across these attacks, harmfulness is no longer only a property of the transcript. It can be introduced, strengthened, or hidden through the acoustic path. A speech-specific threat model should therefore include harmful spoken prompts, distracting non-speech audio, adversarial waveform perturbations, multimodal audio-text jailbreaks, speech-specific jailbreak prompts, and sparse token-level manipulation.

\subsubsection{Output-Side Misuse: Synthetic Voice and Impersonation}

Speech generation introduces a different class of risk. A generated voice can carry both content and an apparent speaker identity. Voice cloning is therefore more than harmful text generation with audio output; it can support impersonation, fraud, reputational harm, evidence manipulation, and social engineering. A speech assistant that can listen, reason, act, and reply in a realistic voice may improve accessibility and interaction, but it can also scale impersonation if voice generation is not constrained.

Voice cloning misuse can be grouped into three categories:

\begin{itemize}
\item \textbf{Identity impersonation:} a malicious actor uses a cloned voice to appear to be another person.
\item \textbf{Evidentiary manipulation:} synthetic audio is presented as proof that someone said something.
\item \textbf{Social engineering:} a cloned voice increases trust in scams, phishing, coercion, or fraudulent instructions.
\end{itemize}

Each category calls for different controls. Identity impersonation requires consent and anti-cloning protection; evidentiary manipulation requires provenance and forensic verification; social engineering requires organizational procedures such as callback verification, transaction limits, and human review.

\subsubsection{Privacy Leakage: Speaker Identity and Paralinguistic Attributes}

Speech privacy extends beyond the literal words being spoken. Even benign speech can reveal speaker identity, accent, age, emotional state, or background context. Speech LLMs that process raw or lightly processed audio may therefore collect more personal information than the user intended to provide. Stored audio and speaker embeddings can also become biometric records that later enable identity recovery or voice cloning.

This distinction separates two related risks. Voice cloning and deepfakes concern synthetic speech as an output misuse problem. Speaker privacy concerns risks that arise from the user's input speech, even when the content is harmless. Raw audio, voice embeddings, transcripts, and paralinguistic features can all contribute to identity leakage. Speech LLM systems should therefore treat speech as both conversational content and personal data.

\subsection{Audio Safety Evaluation and Benchmarking}

As audio jailbreak methods become more varied, benchmark design becomes central to safety evaluation. Early demonstrations often used small prompt sets or single-model case studies. Newer benchmarks aim to test attacks, models, defenses, perturbations, and evaluation protocols under more controlled settings.

\subsubsection{Unified Jailbreak Benchmarks}

\textbf{JALMBench} provides a broad framework for evaluating jailbreak vulnerabilities in large audio-language models. Its dataset includes harmful query samples, text-transferred jailbreak samples, and audio-originated jailbreak samples. It synthesizes audio variants with different accents, genders, TTS methods, and languages, and evaluates multiple LALMs across several attack and defense methods. The paper also reports evaluator reliability analysis, including repeated judging, cross-model consistency, and human consistency verification \citep{peng2026jalmbenchbenchmarkingjailbreakvulnerabilities}. These details are important because benchmark results depend not only on the prompt set, but also on how harmfulness is judged.

The JALMBench results support several design lessons. Non-adversarial harmful queries can have different attack success rates in text and audio modalities. Stronger attacks such as AdvWave can achieve high attack success under the benchmark's settings. Topic sensitivity, attack efficiency, voice diversity, and architecture effects can all influence safety outcomes. These findings suggest that safety evaluation should report accent, speaker, and language variation explicitly, because robustness observed under one speech variety may not generalize to another \citep{peng2026jalmbenchbenchmarkingjailbreakvulnerabilities}.

\subsubsection{Perturbation-Based Benchmarking}

\textbf{Audio Jailbreak} provides a complementary benchmark focused on realistic speech and perturbation. Its base dataset converts adversarial text prompts into speech across policy-violating categories. The paper also introduces an Audio Perturbation Toolkit that applies transformations in the time, frequency, and amplitude domains while aiming to preserve semantic consistency. These transformations test whether safety behavior remains stable under speed changes, noise, clipping, pitch shifts, and related distortions \citep{song2025audiojailbreakopencomprehensive}. This matters because deployed speech systems rarely receive clean, studio-quality audio.

\textbf{Jailbreak-AudioBench} combines a toolbox, dataset, and benchmark for studying audio jailbreaks. It introduces explicit and implicit jailbreak audio examples and evaluates multiple large audio-language models. One contribution is its focus on audio editing as a jailbreak mechanism: harmful semantics can be embedded or transformed through audio operations rather than simply read aloud as a standard prompt \citep{cheng2026jailbreakaudiobenchindepthevaluationanalysis}. This is especially relevant for end-to-end Speech LLMs, which may preserve acoustic information that cascaded ASR-to-LLM systems discard.

\subsubsection{Evaluation Requirements}

Across these benchmark papers, several evaluation requirements emerge:

\begin{itemize}
\item Models should be tested on clean harmful spoken prompts.
\item Models should be tested on perturbed harmful prompts.
\item Benchmarks should include text-transferred and audio-originated jailbreaks.
\item Evaluation should include multimodal audio-text attacks.
\item Speaker, accent, language, and environmental variation should be reported.
\item Defenses should be evaluated together with benign utility.
\end{itemize}

The benchmark gap is therefore not only about dataset size. Audio safety scores depend on perturbation choices, voice diversity, judge stability, target-model assumptions, and benign-utility trade-offs. Benchmark papers should report target models, judges, safety categories, attack assumptions, and whether evaluation relies on human review, model-based judging, or rule-based filtering.

\subsection{Mitigation: Guardrails, Anti-Cloning, and Privacy Defenses}

The attack and benchmark literature shows that text-only moderation is insufficient for Speech LLMs. If an attack is encoded in the acoustic signal, a filter applied only to the transcribed text may miss it. If a system removes too much acoustic information before model processing, it may also harm useful capabilities such as accent robustness, emotion understanding, speaker adaptation, or paralinguistic reasoning. Mitigation therefore has to cover jailbreak guardrails, anti-cloning defenses, privacy protection, and deployment monitoring.

\subsubsection{Audio-Language Guardrails}

\textbf{ALMGuard} provides one example of an audio-language-specific defense. The paper argues that defenses transferred from traditional audio adversarial attacks or text jailbreak defenses are limited against ALM-specific threats. The method identifies safety shortcuts in model activations and introduces safety shortcut activation perturbations as a runtime defense. It also uses a Mel-Gradient Sparse Mask to preserve benign utility. The paper evaluates the defense against multiple attack methods and includes analyses of benign task preservation \citep{jin2025almguardsafetyshortcutsguardrails}. This supports the need for guardrails that operate at the representation or activation level, not only at the transcript level.

Current guardrail research does not yet cover all parts of the threat model equally. Speech-specific jailbreak prompt transformations and sparse token-level manipulation remain only partially addressed by existing defenses. Future defenses may need to combine input detection, token-level robustness analysis, activation-level guardrails, and output monitoring.

\subsubsection{Anti-Cloning Defenses}

For speech synthesis and voice cloning, defenses must also address output-side misuse. \textbf{E2E-VGuard} focuses on unauthorized voice cloning in production LLM-based end-to-end speech synthesis. The paper's threat model describes API-based cloning workflows in which users upload only audio, while the platform internally performs ASR and TTS training. The method protects both timbre and pronunciation: timbre protection uses an encoder ensemble and MFCC-based feature loss to reduce speaker similarity in synthesized audio, while pronunciation protection attacks ASR transcription so that fine-tuning receives incorrect text-audio pairs. The paper also uses a psychoacoustic model to keep perturbations difficult for humans to perceive and evaluates across open-source synthesizers, commercial APIs, ASR systems, languages, and speakers \citep{zhang2025e2evguardadversarialpreventionproduction}. This shows that anti-cloning defense has to consider the full speech synthesis pipeline, not only the final waveform.

\subsubsection{User-Side Voice Privacy Protection}

User-side privacy protection is another defense direction. \textbf{Enkidu} proposes real-time audio privacy protection against personalized voice deepfake threats using universal frequency-domain perturbations. The method does not assume white-box access to the attacker's cloning model. Instead, it uses black-box knowledge, few-shot training, and frequency-domain perturbations designed to preserve perceptual quality while reducing cloneability. The experiments evaluate privacy protection, audio quality, real-time efficiency, and robustness across settings \citep{feng2025enkiduuniversalfrequentialperturbation}. This is relevant for users who want to publish or transmit speech while reducing the risk that their voice can later be cloned.

\subsubsection{Adaptive De-Anonymization and Purification Attacks}

Privacy defenses should be evaluated against adaptive attackers. \textbf{VoxATtack} studies a multimodal attack against voice anonymization systems. Its key insight is that anonymized speech may still leak identity when acoustic and textual information are combined. The method extracts both acoustic embeddings and text embeddings, then fuses them for de-anonymization. Experiments against VoicePrivacy benchmark settings show that multimodal attacks can outperform acoustic-only baselines \citep{aloradi2026voxattackmultimodalattackvoice}. For Speech LLMs, this is important because these systems often process both audio and text representations. Even if the voice is transformed, the transcript may contain linguistic habits, topic choices, names, or discourse patterns that help re-identify the speaker.

\textbf{VocalBridge} further challenges perturbation-based privacy defenses. The paper studies an adaptive attacker who purifies protected speech before using it for voice cloning or speaker verification. Its method uses a diffusion-bridge model in the EnCodec latent space to map protected speech back toward clean speech, with an optional Whisper-guided phoneme conditioning variant. The paper defines user objectives such as imperceptibility, verifiability, and inimitability, and attacker objectives such as restoring verifiability and breaking inimitability. Its evaluation across proactive defenses, voice synthesis tools, and speaker verification systems suggests that perturbation-based protections may be fragile against purification \citep{abbasihafshejani2026vocalbridgelatentdiffusionbridgepurification}.

\subsubsection{Layered Defense Architecture}

A practical mitigation architecture for Speech LLMs should include several layers:

\begin{itemize}
\item \textbf{Input-level defenses:} inspect audio for suspicious transformations, hidden instructions, adversarial perturbations, or abnormal acoustic patterns.
\item \textbf{Representation-level defenses:} monitor acoustic embeddings or hidden states for attack-like behavior.
\item \textbf{Transcription-level defenses:} inspect ASR output while recognizing that transcription is not sufficient on its own.
\item \textbf{Output-level defenses:} moderate generated text and generated speech before delivery.
\item \textbf{Privacy controls:} minimize raw audio storage, separate content processing from identity processing, and document whether speaker embeddings are stored.
\item \textbf{Deployment-level controls:} include logging, rate limits, human escalation, and incident response for high-risk applications.
\end{itemize}

These defenses shift Speech LLM safety from a single-filter problem to a layered systems problem. Because attacks may target the waveform, transcript, internal representation, generated speech, or deployment workflow, defenses should be tested under adaptive conditions.

\subsection{Identity Misuse and Synthetic Speech Provenance}

Synthetic speech risks are related to privacy risks, but they are not the same. The previous subsection focused on defenses against jailbreaks, cloning, and identity leakage. This subsection focuses on generated speech after it leaves the model: whether it can be detected, attributed, and traced.

\subsubsection{Deepfake Detection Under Realistic Conditions}

The \textbf{Perturbed Public Voices (P$^{2}$V)} dataset examines the weakness of audio deepfake detectors under realistic conditions. Its dataset includes identity-consistent transcripts, environmental and adversarial noise, and voice cloning systems from multiple years. The paper evaluates several detectors and shows that performance drops under stronger perturbations and newer cloning methods \citep{gao2025perturbedpublicvoicesp2v}. This supports the argument that deepfake detection should not be validated only on clean or closed-set benchmarks. Speech LLM governance should assume that generated or cloned audio may be compressed, re-recorded, mixed with background noise, or produced by newer synthesis systems than those seen during detector training.

\subsubsection{Robustness Limits of Audio Watermarking}

Watermarking and provenance mechanisms help identify AI-generated speech and support misuse investigation. However, audio can be compressed, resampled, re-recorded, denoised, edited, cropped, transformed by neural codecs, or regenerated by another model. Each operation can weaken or remove a watermark.

A systematic study of audio watermark robustness surveys 22 watermarking schemes and evaluates removal attacks across signal-level, physical-level, and AI-induced distortions. The paper categorizes watermarking approaches, builds an evaluation pipeline, and tests robustness under many removal configurations. The authors conclude that no surveyed scheme is robust against all tested distortions in practice \citep{wen2025sokrobustaudiowatermarking}. This is a caution for Speech LLM deployment: watermarking can support provenance, but it cannot be the only safeguard.

\textbf{Deep Audio Watermarks are Shallow} focuses on limitations of post-hoc speech watermarking. The paper analyzes how post-hoc methods manipulate low-level audio features after generation and shows that transformation-based removal can reduce detectability while preserving audio quality. Its evaluation shows why watermark robustness depends on the attacker's ability to process or transform the signal without knowing the watermarking scheme \citep{oreilly2025deepaudiowatermarksshallow}. This suggests that post-hoc watermarking should be supplemented by generation-time provenance, secure metadata, or model-integrated methods.

\subsubsection{Generation-Time and Dual-Provenance Watermarking}

Several newer methods propose stronger watermarking strategies. \textbf{MelShield} embeds keyed payloads into Mel-spectrogram representations during TTS generation. Its method integrates watermarking into the TTS pipeline rather than relying only on post-processing, and its experiments evaluate attribution accuracy, robustness, and audio quality. This is useful for Speech LLMs because speech generation systems often operate through intermediate acoustic representations such as Mel-spectrograms, making internal watermarking more natural than waveform-only post-processing \citep{jin2026melshieldrobustmeldomainaudio}.

\textbf{DualMark} extends provenance from model-level attribution to both model and training-data attribution. The paper defines a dual-level attribution problem, embeds model and data watermarks into Mel-spectrogram representations during training, and uses watermark consistency loss so generated audio preserves decodable provenance. Its experiments compare model-level and data-level attribution, ablate the watermark consistency loss, and evaluate robustness against pruning, compression, resampling, and noise \citep{yang2025dualmarkidentifyingmodeltraining}. This is relevant because accountability may require identifying both the generating model and the data source.

\textbf{AWARE} proposes watermarking through adversarial resistance to edits. Its method is designed around the possibility that watermarked audio may be edited, cropped, compressed, or desynchronized. Instead of relying only on fixed simulated distortions, it trains the watermarking system with adversarial resistance and uses a detector that aggregates temporal evidence. This direction is important because synthetic speech may be transformed many times before detection \citep{pavlović2026awareaudiowatermarkingadversarial}.

The watermarking literature reframes provenance as a robustness-under-transformation problem. The question is not only whether a watermark can be detected on clean generated audio, but whether it survives compression, re-recording, editing, adversarial removal, and model-based transformation.

\subsubsection{Traceability Stack}

For Speech LLM deployment, provenance should be treated as a stack rather than a single watermark. A traceability system may include:

\begin{itemize}
\item Audio watermarking.
\item Secure metadata.
\item Model version records.
\item Consent records.
\item User identity verification for voice cloning.
\item Audit logs and incident response procedures.
\end{itemize}

In public-facing or high-risk settings, generated speech should remain traceable even when the audio watermark fails. The secure-metadata layer is now being standardized outside the research literature: the C2PA Content Credentials specification attaches cryptographically signed provenance manifests to media, and news organizations have begun deploying it against deepfakes and disinformation \citep{strickland2024credentials}. Watermarking is useful but insufficient: it supports accountability, but it has to be combined with such signed metadata, policy, monitoring, and operational safeguards.

\subsection{Governance, Documentation, and Compliance Gaps}

The traceability stack above describes technical provenance mechanisms. This subsection turns to documentation and governance. Speech LLMs need documentation that covers not only model accuracy, but also intended use, prohibited use, data provenance, speaker privacy, audio safety testing, watermarking, logging, and incident response. General AI governance frameworks provide useful starting points, but they need speech-specific extensions. Recent policy scholarship on generative-AI governance emphasizes that effective regimes must span the full model lifecycle---development, deployment, and post-deployment monitoring---rather than regulating a single release point \citep{taeihagh2025governance}, a lifecycle view that maps directly onto the speech-specific requirements below.

\subsubsection{Model Documentation}

\textbf{Model Cards} offer a starting point for standardized model reporting. The paper proposes short documents accompanying trained models, including intended use, performance evaluation, ethical considerations, and performance across relevant groups and conditions. Although the original framework focuses on machine learning models more broadly, the same structure can be adapted to Speech LLMs \citep{mitchell2019modelcards}. A Speech LLM model card should report supported languages and accents, evaluated noise conditions, speech input and output capabilities, audio jailbreak testing, deepfake risk controls, privacy protections, latency constraints, and whether raw audio or speaker embeddings are stored.

\subsubsection{Machine-Readable Governance}

Machine-readable documentation may help connect policy requirements to system behavior. \textbf{AI Cards} propose human- and machine-readable documentation for AI systems, including technical specifications, context of use, and risk management information inspired by EU AI Act requirements. This type of structured documentation is useful for Speech LLMs because deployment risk depends heavily on context: a speech model used for entertainment has different requirements from one used in healthcare, finance, education, or public services \citep{golpayegani2024aicardsappliedframework}.

\textbf{Policy Cards} further extend this idea to runtime governance. They encode operational, regulatory, and ethical constraints as machine-readable artifacts for autonomous AI agents and emphasize auditability, compliance, transparency, and runtime policy enforcement \citep{mavracic2025policycards}. For speech assistants that can call tools or act in real time, runtime governance can help translate high-level rules into operational controls.

\subsubsection{Governance Requirements for Speech LLMs}

A governance framework for Speech LLMs should include at least five components:

\begin{itemize}
\item \textbf{Data governance:} require consent and provenance documentation for speech data, especially when voices are identifiable or cloneable.
\item \textbf{Safety governance:} require audio-specific red teaming, including clean harmful speech, perturbed audio, hidden-semantics attacks, and multimodal audio-text jailbreaks.
\item \textbf{Privacy governance:} restrict raw audio retention, speaker embedding storage, and sensitive inference from paralinguistic cues.
\item \textbf{Provenance governance:} require synthetic speech disclosure, watermarking where feasible, and secure audit trails.
\item \textbf{Operational governance:} define escalation paths, human review, incident reporting, and post-deployment monitoring.
\end{itemize}

Most existing documentation frameworks are modality-general, while Speech LLMs create speech-specific risks. Voice is an interaction medium, a biometric signal, and a synthetic media carrier. Governance documentation should therefore specify how the system handles audio data, speaker identity, generated speech, and misuse reports.

\subsection{Design Implications}

The reviewed literature suggests several design implications for Speech LLM development and deployment.

\begin{itemize}
\item \textbf{Data collection:} systems should document consent, voice provenance, and speaker privacy risks.
\item \textbf{Training and adaptation:} developers should examine whether multimodal alignment changes safety behavior inherited from the text backbone.
\item \textbf{Evaluation:} benchmarks should include audio jailbreaks, perturbation robustness, accent and language variation, privacy leakage, and deepfake misuse.
\item \textbf{Deployment:} systems should use layered guardrails, provenance mechanisms, logging, and human escalation.
\item \textbf{Post-deployment monitoring:} organizations should track misuse patterns, update defenses, and disclose incidents where appropriate.
\end{itemize}

The central design point is that speech is not just text with sound. It is a high-dimensional signal that can carry semantic content, speaker identity, emotion, and contextual information. Audio jailbreaks attack speech as a command channel. Adversarial perturbations attack speech as a signal. Voice cloning attacks speech as identity. Deepfakes attack speech as evidence. Privacy attacks exploit speech as personal data. Watermarking and provenance govern speech as synthetic media. Speech LLM safety therefore depends on technical robustness, privacy protection, provenance, and operational controls. Section~\ref{sec:future_directions} discusses open problems in audio-specific red teaming, adaptive defenses, privacy-preserving speech representations, and provenance under audio transformation.
